\documentclass[12pt,a4paper]{article}
\usepackage[utf8]{inputenc}
\usepackage[T1]{fontenc}
\usepackage[margin=1in]{geometry}
\usepackage{amsmath}
\usepackage{graphicx}
\usepackage{textcomp}
\usepackage{setspace}
\usepackage{microtype}
\usepackage{threeparttable}
\usepackage{booktabs}
\usepackage{adjustbox}
\usepackage{array}
\usepackage{multirow}
\usepackage{setspace}
\usepackage{underscore}
\usepackage{enumitem}
\usepackage{float}
\usepackage{pdflscape}
\usepackage{caption}
\usepackage[authoryear]{natbib}
\newcommand{\sym}[1]{\ensuremath{^{#1}}}

\usepackage{color}
\usepackage[dvipsnames]{xcolor}
\definecolor{darkblue}{RGB}{0.,0.,139.}

\usepackage{hyperref}
\usepackage{cleveref} 
\hypersetup{unicode=true,bookmarksnumbered=true,bookmarksopen=true,bookmarksopenlevel=3,
  breaklinks=true,pdfborder={0 0 0},colorlinks,citecolor=darkblue,filecolor=darkblue,
  linkcolor=darkblue,urlcolor=darkblue}

\onehalfspacing

\newcommand{\tablepath}{tables/}

\title{Career Incentives and Research Topic Choice Among Economists}
\author{Sumedha Ray}
\date{\today}

\begin{document}

\maketitle
{
\singlespacing
\begin{abstract}
Do economists who study sensitive topics such as race, religion, discrimination, and gender, face career penalties, and do they respond by strategically shifting their research agendas around tenure? This paper constructs a novel dataset linking dissertation topic choices from the JEL Annual Dissertations Lists (1990--2024) to post-PhD career outcomes from IDEAS/RePEC. Preliminary descriptive evidence from the 2024 cohort shows that women are substantially more likely than men to write dissertations on sensitive topics, particularly at top PhD programs.
\end{abstract}
}

\newpage

\section{Introduction}
\label{sec:intro}
\begin{itemize}
    \item Academic economists face strong incentives to publish in top journals and secure tenure.
    \item Research on sensitive topics (race, religion, discrimination, gender) may be harder to publish or perceived as less rigorous, creating career risk.
    \item This paper asks: (1) Are economists who study sensitive topics penalized in their careers? (2) Do they strategically shift research agendas around tenure?
    \item I use the JEL Annual Dissertations Lists to track dissertation topic choice from 1990--2024 and link to career outcomes via IDEAS/RePEC.
    \item Preliminary evidence: women are substantially more likely to write sensitive-topic dissertations, especially at top programs.
\end{itemize}

\section{Literature Review}
\label{sec:litreview}
\begin{itemize}
    \item \citet{ginther2004}: women less likely to receive tenure in economics even conditional on productivity.
    \item \citet{bayer2016}: overview of underrepresentation of women and minorities; structural barriers in the profession.
    \item \citet{card2020}: evidence of gender bias in peer review at top journals.
    \item \citet{hengel2022}: papers by women take longer through peer review and are held to higher standards.
    \item \citet{doleac2022}: documents who researches race and policing; notes underrepresentation of minority scholars.
    \item \citet{koffi2021}: women more likely to pursue novel and unconventional research, which may carry higher career variance.
    \item \citet{akerlof2020}: career incentives can distort research choices and perpetuate existing paradigms.
\end{itemize}

\section{Data}
\label{sec:data}
\begin{itemize}
    \item \textbf{JEL Annual Dissertations Lists}: published annually in the \textit{Journal of Economic Literature}, available on JSTOR. Contains dissertation title, author name, institution, graduation year, and JEL codes. Sample: 1990--2024.
    \item \textbf{Sensitive topic classification}: dissertation flagged as sensitive if title contains keywords related to race, ethnicity, religion, gender, or discrimination, OR if JEL codes include categories such as J15, J16, Z12. Formally:
    \begin{equation}
        S_i = \mathbf{1}\left[\text{title}_i \in \mathcal{K} \;\cup\; \text{JEL}_i \in \mathcal{J}\right]
    \end{equation}
    where $\mathcal{K}$ is the keyword set and $\mathcal{J}$ is the set of sensitive JEL codes.
    \item \textbf{Career outcomes}: IDEAS/RePEC provides publication records, citation counts, and institutional affiliations. Merged to dissertation records using author name, institution, and graduation year with fuzzy matching.
    \item \textbf{Current status}: 2024 cohort parsed and analyzed. Full 1990--2024 merge in progress.
\end{itemize}

\section{Empirical Strategy}
\label{sec:empirical}
\begin{itemize}
    \item \textbf{Cross-sectional}: estimate gender gap in sensitive-topic dissertation choice, with cohort and program fixed effects:
    \begin{equation}
        S_{ic} = \alpha + \beta\,\text{Female}_{ic} + \gamma\,\text{TopProgram}_{ic} + \delta\,(\text{Female} \times \text{TopProgram})_{ic} + \theta_c + \varepsilon_{ic}
    \end{equation}
    \item \textbf{Career outcomes}: regress publications, citations, and placement rank on sensitive-topic indicator, controlling for gender, program rank, cohort, and field:
    \begin{equation}
        Y_{ic} = \alpha + \beta\, S_{ic} + \mathbf{X}_{ic}'\gamma + \theta_c + \lambda_p + \varepsilon_{ic}
    \end{equation}
    \item \textbf{Research agenda shift}: event-study around expected tenure year ($\approx$6--8 years post-PhD), testing whether sensitive-topic researchers reduce their share of sensitive publications approaching tenure.
    \item Key identification challenge: separating genuine preference change from strategic self-censorship.
\end{itemize}

\section{Results}
\label{sec:results}
\begin{itemize}
    \item Women are substantially more likely than men to write sensitive-topic dissertations in the 2024 cohort (see Table~\ref{tab:summary}).
    \item This gender gap is largest at top-10 and top-20 programs (see Figure~\ref{fig:gender_gap}).
    \item Gap persists across all sensitive subcategories; largest for gender-related topics.
    \item Sensitive-topic dissertations concentrated in labor economics and development economics.
    \item \textit{Longitudinal career outcomes results pending completion of IDEAS/RePEC merge.}
\end{itemize}

\section{Conclusion}
\label{sec:conclusion}
\begin{itemize}
    \item Preliminary evidence points to a large gender gap in sensitive-topic dissertation research, particularly at top programs.
    \item Next steps: complete IDEAS/RePEC merge for 1990--2024 cohorts; estimate career penalty and research-agenda-shift specifications.
    \item If sensitive-topic research is penalized, the profession may be systematically underproducing knowledge on some of the most policy-relevant questions.
\end{itemize}

\vfill
\pagebreak{}
\cleardoublepage
\begin{spacing}{1.0}
\bibliographystyle{jpe}
%\nocite{*}
\bibliography{myfile}
\addcontentsline{toc}{section}{References}
\end{spacing}
\cleardoublepage

%--------------------------------------------------
% Tables and Figures
%--------------------------------------------------

\section*{Figures and Tables}
\label{sec:figures-tables}

\begin{table}[h!]
\centering
\caption{Summary Statistics: 2024 JEL Annual Dissertations List}
\label{tab:summary}
\begin{tabular}{lcc}
\toprule
 & Men & Women \\
\midrule
Total dissertations & [N] & [N] \\
Share sensitive topic (\%) & [x] & [x] \\
Top-10 program (\%) & [x] & [x] \\
\bottomrule
\end{tabular}
\caption*{\textit{Notes:} Data from 2024 JEL Annual Dissertations List. To be populated with actual values.}
\end{table}

\newpage
\begin{figure}[h!]
\centering
% \includegraphics[width=0.85\textwidth]{figures/gender_gap_by_tier.pdf}
\caption{Gender Gap in Sensitive-Topic Dissertations by Program Tier, 2024}
\label{fig:gender_gap}
\caption*{\textit{Notes:} Share of dissertations classified as sensitive, by gender and program tier. Figure to be inserted once R output is finalized.}
\end{figure}

\end{document}